%%%%%%%%%%%%%%%%%%%%%%%%%%%%%%%%%%%%%%%%%%%%%%%%%%%%%%%%%%%%%%%%%%%%%%
% Lớp 12 | Nguyên hàm, tích phân | Tích phân
% Dạng 2D4?2-5: Phương pháp đổi biến số cơ bản
% Số câu: 16  (sinh tự động bởi tach_cau.py; từ nay thầy sửa trực tiếp file này)
%%%%%%%%%%%%%%%%%%%%%%%%%%%%%%%%%%%%%%%%%%%%%%%%%%%%%%%%%%%%%%%%%%%%%%

%%%%%%%%%%%%%%%% TRẮC NGHIỆM 4 PHƯƠNG ÁN %%%%%%%%%%%%%%%%

% nguồn: Phan1_Lop_12_Dai_so_Chuong_4
\begin{ex}%[2D4H2-5]%[ma:6e79f204]
	Tích phân $\displaystyle\int_0^1{{\mathrm e^{3x}}\mathrm{\,d}x}$  bằng
	 \choice
		 {${\mathrm e^3}+\dfrac12$}
		{$ \mathrm e-1$} 
		 {\True $\dfrac{{\mathrm e^3}- 1}3$}
		{${\mathrm e^3}-  1$}
	 \loigiai{
		Ta có  \[\displaystyle\int_0^1{{\mathrm e^{3x}}\mathrm{\,d}x}=\dfrac13\displaystyle\int_0^1{{\mathrm e^{3x}} \mathrm{\,d}\left(3x\right)}=\dfrac13{\mathrm e^{3x}} \Bigg|_0^1=\dfrac{{\mathrm e^3}-1}3.\]
	} 
\end{ex}

% nguồn: Phan1_Lop_12_Dai_so_Chuong_4
\begin{ex}%[2D4H2-5]%[ma:90b5fe81]
	Cho  hàm số $f(x)=\heva{&x+1&&\text{ khi }x\geq 0\\
	&\mathrm{e}^{2x}&&\text{  khi }x<0}$. Tích phân $I=\int_{-1}^2f(x)\mathrm{\,d}x$ có giá trị bằng bao nhiêu?
	 \choice
		{$I=\dfrac{3\mathrm{e}^2-1}{\mathrm{e}^2}$} 
		{\True  $I=\dfrac{9\mathrm{e}^2-1}{2\mathrm{e}^2}$}
		 {$I=\dfrac{11\mathrm{e}^2-11}{2\mathrm{e}^2}$}
		 {$I=\dfrac{7\mathrm{e}^2+1}{2\mathrm{e}^2}$}
	\loigiai{Ta có \allowdisplaybreaks\begin{eqnarray*}
		I&=&\int_{-1}^2f(x)\mathrm{\,d}x=\int_{-1}^0f(x)\mathrm{\,d}x+\int_0^2f(x)\mathrm{\,d}x\\ 
		&=&\int_{-1}^0\mathrm{e}^{2x}\mathrm{\,d}x+\int_0^2(x+1)\mathrm{\,d}x=\left.\dfrac{1}{2}\mathrm{e}^{2x}\right|_{-1}^0+\left.\left(\dfrac{x^2}{2}+x\right)\right|_0^2\\ 
		&=&\dfrac{1}{2}\left(1-\mathrm{e}^{-2}\right)+\left(\dfrac{4}{2}+2\right)=\dfrac{9}{2}-\dfrac{1}{2\mathrm{e}^2}=\dfrac{9\mathrm{e}^2-1}{2\mathrm{e}^2}. 
	\end{eqnarray*}
 }
\end{ex}

% nguồn: Phan1_Lop_12_Dai_so_Chuong_4
\begin{ex}%[2D4H2-5]%[ma:d12137df]
	Nếu $\displaystyle\int_0^{\pi }{f(x)\,\mathrm{\,d}x}=3$ thì $\displaystyle\int_0^{\pi }{\left[f(x)+\sin\dfrac x2\right]\,\mathrm{\,d}x}$ bằng 
	 \choice
		{$10$} 
		 {$6$}
		 {$12$}
		{\True $5$} 
	\loigiai{ 
		Ta có \[\displaystyle\int_0^{\pi }{\left[f(x)+\sin\dfrac x2\right]\,\mathrm{\,d}x}=\displaystyle\int_0^{\pi }{f(x)\,\mathrm{\,d}x}+\displaystyle\int_0^{\pi }{\sin\dfrac  x2\,\mathrm{\,d}x}=3- 2\cos\dfrac x2 \Bigg|_0^{\pi }=3-2\left(0-1\right)=5.\]
	 }
\end{ex}

% nguồn: Phan1_Lop_12_Dai_so_Chuong_4
\begin{ex}%[2D4V2-5]%[ma:4285f163]
	Cho hàm số $y=f(x)$ có đạo hàm trên $\mathbb{R}$, thỏa mãn  $f(x)=\heva{&2x&&\text{ khi }x\geq 0\\&x^3+x+1&&\text{ khi }x<0}$. Tính $\displaystyle\int_0^2f(x-1)\mathrm{\,d}x$.
	\choice 
		 {$\dfrac{3}{2}$}
		 {$\dfrac{3}{4}$}
		{\True $\dfrac{5}{4}$} 
		{$\dfrac{11}{4}$} 
	\loigiai{ 
		Ta có \allowdisplaybreaks\begin{eqnarray*}
		 \int_0^2f(x-1)\mathrm{\,d}x&=&\int_{-1}^1f(x)\mathrm{\,d}x\\
		&=&\int_{-1}^0f(x)\mathrm{\,d}x+\int_0^1f(x)\mathrm{\,d}x\\ 
		&=&\int_{-1}^0\left(x^3+x+1\right)\mathrm{\,d}x+\int_0^12x\mathrm{\,d}x\\ 
		&=&\left(\dfrac{x^4}{4}+\dfrac{x^2}{2}+x\right)\bigg|_{-1}^0+x^2\bigg|_0^1\\ 
		&=&-\dfrac{1}{4}-\dfrac{1}{2}+1+1\\ 
		&=&\dfrac{5}{4}. 
	\end{eqnarray*}
} 
\end{ex}

% nguồn: Phan1_Lop_12_Dai_so_Chuong_4
\begin{ex}%[2D4V2-5]%[ma:9ee62a41]
	Cho hàm số $f(x)=\heva{&x^2-x+1&&\text{ khi }x\geq \dfrac{1}{2}\\&2x+\dfrac{7}{4}&&\text{  khi }x<\dfrac{1}{2}}$. Tính tích phân $\displaystyle\int_0^{\tfrac{\pi}{2}}f\left(\sin ^2x\right)\sin 2x\mathrm{\,d}x$.
	 \choice
		{\True  $\dfrac{37}{24}$}
		 {$-\dfrac{37}{24}$}
		 {$\dfrac{24}{37}$}
		{$\dfrac{17}{24}$} 
	\loigiai{ 
		Ta có \allowdisplaybreaks\begin{eqnarray*}
		 \int_0^{\tfrac{\pi}{2}}f\left(\sin ^2x\right)\sin 2x\mathrm{\,d}x&=&\int_0^{\tfrac{\pi}{2}}f\left(\sin ^2x\right)\cdot 2\sin x\cos x\mathrm{\,d}x\\
		&=&\int_0^{\tfrac{\pi}{2}}2\sin x.f\left(\sin ^2x\right)\mathrm{d}(\sin  x)\\
		&=&\int_0^{\tfrac{\pi}{2}}f\left(\sin  ^2x\right)\mathrm{d}\left(\sin ^2x\right)\\
		 &=&\int_0^1f(u)du\\
		&=&\int_0^1f(x)\mathrm{\,d}x\\ 
		 &=&\int_0^{\tfrac{1}{2}}\left(2x+\dfrac{7}{4}\right)\mathrm{\,d}x+\int_{\tfrac{1}{2}}^1\left(x^2-x+1\right)\mathrm{\,d}x\\
		 &=&\dfrac{37}{24}.
	\end{eqnarray*}
 }
\end{ex}

% nguồn: Phan1_Lop_12_Dai_so_Chuong_4
\begin{ex}%[2D4V2-5]%[ma:3c4076d7]
	Biết $I=\displaystyle\int_0^{\tfrac{\pi }2}{\dfrac{x+x\cos x-{{\sin }^3}x}{1+\cos x}\mathrm{\,d}x}=\dfrac{{{\pi }^2}}a-\dfrac bc$. Trong đó $a$, $b$, $c$ là các số nguyên dương, phân số $\dfrac bc$  tối giản. Tính $T={a^2}+{b^2}+{c^2}$.
	 \choice
		{$T=50$} 
		 {$T=59$}
		{$T=16$} 
		 {\True $T=69$}
	\loigiai{ 
		 Ta có
		\begin{align*}
			 I&=\displaystyle\int_0^{\tfrac{\pi }2}{\dfrac{x+x\cos x-{{\sin }^3}x}{1+\cos x}\mathrm{\,d}x}\\
			&=\displaystyle\int_0^{\tfrac\pi2}{\left(x-\dfrac12\sin2x\right)}\mathrm{\,d}x\\ 
			 &=\left(\dfrac{x^2}2+\dfrac14\cos2x\right)\Bigg|_0^{\tfrac\pi2}\\
			&=\dfrac{{{\pi  }^2}}8-\dfrac12.
		\end{align*}
		 Khi đó $\heva{&a=8\\&b=1\\&c=2}
		\Rightarrow{a^2}+{b^2}+{c^2}=69$. 
	} 
\end{ex}

% nguồn: Phan1_Lop_12_Dai_so_Chuong_4
\begin{ex}%[2D4V2-5]%[ma:05e83ddc]
	Có  bao nhiêu số thực $b$ thuộc khoảng $\left(\pi\,;\,3\pi \right)$ sao cho $\displaystyle\int_{\pi }^b{4\cos2x\mathrm{\,d}x}=1$?
	\choice 
		{\True  $4$}
		 {$6$}
		 {$8$}
		 {$2$}
	\loigiai{Ta có 
		\begin{align*}
			\displaystyle\int_{\pi }^b{4\cos2x\mathrm{\,d}x}=1&\Leftrightarrow 2\sin2x \Bigg|_{\pi }^b=1\\ 
			 &\Leftrightarrow\sin2b=\dfrac12\\
			&\Leftrightarrow\hoac{ 
			&2b=\dfrac{\pi  }6+k2\pi \\
			&2b=\dfrac{5\pi  }6+k2\pi \\
			}\\ 
			 &\Leftrightarrow\hoac{
			&b=\dfrac{\pi }{12}+k\pi \\ 
			&b=\dfrac{5\pi }{12}+k\pi \\ 
			},\  \left(k\in\mathbb{Z}\right).
		\end{align*}
		Với $ b=\dfrac{\pi }{12}+k\pi $  mà $ b\in\left(\pi;\,3\pi \right)\Rightarrow 1<\dfrac1{12}+k<3$, $\left(k\in\mathbb{Z}\right)\Rightarrow k=1;2\Rightarrow b=\dfrac{13\pi }{12};\dfrac{25\pi }{12}$.\\
		Với $ b=\dfrac{5\pi }{12}+k\pi $ mà $  b\in\left(\pi;\,3\pi \right)\Rightarrow 1<\dfrac5{12}+k<3$, $\left(k\in\mathbb{Z}\right)\Rightarrow k=1;2\Rightarrow b=\dfrac{17\pi }{12};\dfrac{29\pi }{12}$.\\
	Vậy có $4$ số thực $b$  thỏa mãn yêu cầu bài toán.}
\end{ex}

% nguồn: Phan1_Lop_12_Dai_so_Chuong_4
\begin{ex}%[2D4V2-5]%[ma:37a9ed24]
	Tính tích phân $\displaystyle\int_0^{\tfrac\pi2}{{\mathrm{e}^{\cos x}}\sin x\mathrm{\,d}x}$  bằng
	 \choice
		{$1-\mathrm  e$}
		{$\mathrm  e+1$}
		{$\mathrm e$} 
		{\True  $\mathrm e-1$}
	 \loigiai{
		Ta  có
		\[\displaystyle\int_0^{\tfrac\pi2}{{\mathrm{e}^{\cos x}}\sin x\mathrm{\,d}x}=-\displaystyle\int_0^{\tfrac\pi2}{{\mathrm{e}^{\cos x}} \mathrm{\,d}\left(\cos x\right)}=-{\mathrm{e}^{\cos x}} \Bigg|_0^{\tfrac\pi2}=-\left({\mathrm{e}^{\cos\tfrac\pi2}}-{\mathrm{e}^{\cos 0}}\right)=-\left(1-\mathrm e\right)=\mathrm e-1.\] 
	 }
\end{ex}

% nguồn: Phan1_Lop_12_Dai_so_Chuong_4
\begin{ex}%[2D4V2-5]%[ma:581e8132]
	 Cho $I=\displaystyle\int_0^1{\dfrac{\mathrm{\,d}x}{\sqrt{2x+a}}}$, với $ a>0$. Tìm $a$ nguyên để $I\ge1$.
	 \choice
		{\True Không có giá trị nào của  $a$}
		{$ a=0$} 
		 {Vô số giá trị của $a$}
		{$  a=1$}
	 \loigiai{
		Ta  có \[I=\displaystyle\int_0^1{\dfrac{\mathrm{\,d}x}{\sqrt{2x+a}}}=\sqrt{2x+a} \Bigg|_0^1=\sqrt{2+a}-\sqrt a.\]
		Ta có \[I\ge1\Leftrightarrow\sqrt{2+a}\ge1+\sqrt a\Leftrightarrow 2+a\ge a+1+2\sqrt a\Leftrightarrow\sqrt a\le\dfrac12\Leftrightarrow 0\le  a\le\dfrac14.\]
		Với $a>0$, $a$ nguyên thì không có giá trị nào của $a$  thoả mãn.
	 }
\end{ex}

% nguồn: Phan1_Lop_12_Dai_so_Chuong_4
\begin{ex}%[2D4V2-5]%[ma:1bfa1c72]
	Biết $I=\displaystyle\int_1^2{\dfrac12}\left(\dfrac1x-\dfrac1{x+2}\right)\mathrm{\,d}x=a\ln 2+b\ln 3$ với $ a,b\in\mathbb{R}$. Tính  $T={a^2}+{b^3}$.
	 \choice
		{$T=\dfrac18$} 
		 {$T=\dfrac83$}
		 {$T=\dfrac12$}
		 {\True $T=\dfrac38$}
	\loigiai{ 
		Ta  có
		\begin{align*}
			I&=\displaystyle\int_1^2{\dfrac12}\left(\dfrac1x-\dfrac1{x+2}\right)\mathrm{\,d}x\\ 
			&=\dfrac12\left[\ln x-\ln\left(x+2\right)\right]\Bigg|_1^2\\ 
			 &=\dfrac12\ln\dfrac x{x+2}\Bigg|_1^2\\
			&=\dfrac12\left(\ln\dfrac12-\ln\dfrac13\right)\\ 
			&=-\dfrac12\ln  2+\dfrac12\ln 3.
		\end{align*}
		Từ đó  $a=-\dfrac12$, $b=\dfrac12$, vậy $T={a^2}+{b^3}=\dfrac38$.
	} 
\end{ex}

% nguồn: Phan1_Lop_12_Dai_so_Chuong_4
\begin{ex}%[2D4V2-5]%[ma:9be206bc]
	Gọi $a$, $b$ là các số nguyên sao cho $\displaystyle\int_0^2{\sqrt{{\mathrm{e}^{x+2}}}\mathrm{\,d}x}=2a{\mathrm{e}^2}+b\mathrm e$.  Giá trị của ${a^2}+{b^2}$ bằng
	\choice 
		 {$3$}
		 {$8$}
		 {$4$}
		{\True  $5$}
	\loigiai{ 
		 Ta có
		 \[\displaystyle\int_0^2{\sqrt{{\mathrm{e}^{x+2}}}\mathrm{\,d}x}=\displaystyle\int_0^2{{{\left({\mathrm{e}^{x+2}}\right)}^{\tfrac12}}\mathrm{\,d}x}=\displaystyle\int_0^2{{\mathrm{e}^{\tfrac12x+1}}\mathrm{\,d}x}= 2{\mathrm{e}^{\tfrac12x+1}} \Bigg|_0^2=2\mathrm{e}^2-2\mathrm e\Rightarrow\heva{
		&a=1\\ 
		&b=-2.}\] 
	Vậy  ${a^2}+{b^2}=5$.}
\end{ex}

% nguồn: Phan1_Lop_12_Dai_so_Chuong_4
\begin{ex}%[2D4V2-5]%[ma:117f2fd0]
	Cho $\displaystyle\int_0^1{\dfrac{\mathrm{\,d}x}{\sqrt{x+2}+\sqrt{x+1}}}=a\sqrt b-\dfrac83\sqrt a+\dfrac23$, $\left(a,b\in {{\mathbb{N}}^{*}}\right)$. Tính  $ a+2b$
	\choice 
		 {$ \text7$}
		{\True $  \text8$}
		 {$-1$}
		 {$5$}
	\loigiai{ 
		Ta  có
		 \allowdisplaybreaks
		\begin{eqnarray*}
			\displaystyle\int_0^1{\dfrac{\mathrm{\,d}x}{\sqrt{x+2}+\sqrt{x+1}}}&=&\displaystyle\int_0^1{\left(\sqrt{x+2}-\sqrt{x+1}\right)}\mathrm{\,d}x\\ 
			&=&\dfrac23\left(\sqrt{{{\left(x+2\right)}^3}}-\sqrt{{{\left(x+1\right)}^3}}\right)  \Bigg|_0^1\\
			 &=&2\sqrt3-\dfrac83\sqrt2+\dfrac23.
		\end{eqnarray*}
	 Do đó $ a=2$, $b=3$ nên $a+2b=8$.}
\end{ex}


%%%%%%%%%%%%%%%% ĐÚNG/SAI %%%%%%%%%%%%%%%%

% nguồn: Phan1_Lop_12_Dai_so_Chuong_4
\begin{ex}%[2D4H2-5]%[ma:d237d996]
	Cho $f(x)=2x-\sin x$ và $F(x)$ là một nguyên hàm của $f(x)$. Khi  đó
	 \choiceTF
		{\True Cho $G(x)=F(x)+C$,  $C$ là hằng số thì $G(x)$ là một nguyên hàm của $f(x)$}
		{$\displaystyle\int f(x)\mathrm{\,d}x=x^2-\cos x+C$, với  $C$ là hằng số}
		 {\True $\displaystyle\int_{0}^{\tfrac{\pi}{3}}f(x)\mathrm{\,d}x=\dfrac{\pi^2}{9}-\dfrac{1}{2}$}
		{Diện tích hình phẳng giới hạn bởi đồ thị hàm số $f(x)=2x-\sin  x$, trục hoành và hai đường thẳng $x=-1$, $x=2$ bằng $\displaystyle\int_{-1}^{2}(2x-\sin x)\mathrm{\,d}x$}
	 \loigiai{
		\begin{itemchoice}
			\itemch Ta có $\left(G(x)\right)'=\left(F(x)+C\right)'=F'(x)=f(x)$. Do đó $G(x)$ là  một nguyên hàm của $f(x)$.
			\itemch $\displaystyle\int f(x)\mathrm{\,d}x=\displaystyle\int  2x-\sin x\mathrm{\,d}x=x^2+\cos x+C$, với $C$ là hằng số.
			\itemch $\displaystyle\int_{0}^{\tfrac{\pi}{3}}f(x)\mathrm{\,d}x=\left(x^2+\cos x\right)\bigg|_{0}^{\frac{\pi}{3}}=\dfrac{\pi^2}{9}-\dfrac{1}{2}$. 
			 \itemch $S=\displaystyle\int_{-1}^{2}\left|2x-\sin x\right|\mathrm{\,d}x$.
		\end{itemchoice}
	 }
\end{ex}

% nguồn: Phan1_Lop_12_Dai_so_Chuong_4
\begin{ex}%[2D4V2-5]%[ma:9cdfe187]
	Trên khoảng $(0;+\infty)$, cho hàm số  $f(x)=\dfrac{1-3\sqrt{x}}{x}$. Khi đó
	\choiceTF 
		{\True  $\displaystyle\int_1^4f'(x) \mathrm{\,d}x=\dfrac{3}{4}$}
		 {\True $\displaystyle\int_1^4f(x) \mathrm{\,d}x=-6+\ln 4$}
		{$\displaystyle\int_1^4xf(x)  \mathrm{\,d}x=11$}
		 {$\displaystyle\int_1^4\left[f'(x)+2x\cdot f(x)\right] \mathrm{\,d}x=\dfrac{91}{4}$}
	\loigiai{ 
		\begin{itemchoice}
			\itemch {\bf  Đúng}. \\
			$\displaystyle\int_1^4f'(x)  \mathrm{\,d}x
			 =\left.f(x)\right|_1^4=\dfrac{1-3\sqrt{4}}{4}-\dfrac{1-3\sqrt{1}}{1}
			 =\dfrac{3}{4}$.
			\itemch {\bf  Đúng}.
			\begin{eqnarray*}
				\displaystyle\int_1^4f(x) \mathrm{\,d}x&=&\displaystyle\int_1^4\left(\dfrac{1}{x}-3x^{-\tfrac{1}{2}} \right)  \mathrm{\,d}x\\
				&=&\left.\left(\ln x-6\sqrt{x} \right)\right|_1^4\\ 
				&=&(\ln4-6\cdot  2)-(\ln1-6)\\
				&=&\ln4-6. 
			\end{eqnarray*}
			\itemch  {\bf Sai}.\\
			 $\displaystyle\int_1^4f(x) \mathrm{\,d}x
			=\displaystyle\int_1^4\left(1-3x^{\tfrac{1}{2}} \right)  \mathrm{\,d}x
			=\left.\left(x-2x^{\tfrac{3}{2}} \right)\right|_1^4 
			=(4-16)-(1-2)=-11$. 
			\itemch {\bf  Sai}. \\
			$\displaystyle\int_1^4\left[f'(x)+2x\cdot f(x)\right]\mathrm{\,d}x 
			=\displaystyle\int_1^4f'(x)\mathrm{\,d}x+2\cdot  \displaystyle\int_1^4xf(x)\mathrm{\,d}x
			=\dfrac{3}{4}+2\cdot  (-11)
			=-\dfrac{85}{4}$. 
		\end{itemchoice}
	 }
\end{ex}


%%%%%%%%%%%%%%%% TRẢ LỜI NGẮN %%%%%%%%%%%%%%%%

% nguồn: Phan2_Lop_12_Dai_so_Chuong_4
\begin{bt}%[2D4H2-5]%[ma:1ee9d442]
	Cho hàm số $y=f(x)=\heva{&\mathrm{e}^x+1 &\text{ khi }x\ge 0 \\&x^2-2x+2 &\text{ khi }x<0}$. Biết giá trị của \break $\displaystyle I=\int_{\tfrac{1}{\mathrm{e}}}^{\mathrm{e}^2}\dfrac{f\left(\ln x-1\right)}{x}\mathrm{\,d}x =\dfrac{a}{b}+c\mathrm{e}$ với $a; b; c\in \mathbb{Z}$, $b\ne 0$ và $\left(a; b\right)=1$. Giá  trị của $a+b+c$ bằng
	 \shortans{36}
	 \loigiai{
		Xét $\displaystyle I=\int_{\tfrac{1}{\mathrm{e}}}^{\mathrm{e}^2}\dfrac{f(\ln  x-1)}{x}\mathrm{\,d}x$. \\
		Đặt $u=\ln x-1\Rightarrow \mathrm{\,d}u=\dfrac{1}{x}\mathrm{\,d}x$. 
		Đổi cận $\heva{&x=\dfrac{1}{\mathrm{e}}\Rightarrow u=-2\\&x=\mathrm{e}^2\Rightarrow  u=1.}$\\
		Khi  đó
		\begin{eqnarray*}
			\displaystyle  I&=&\int_{-2}^1 f(u)\mathrm{\,d}u =\int_{-2}^1f(x)\mathrm{\,d}x \\ &=&\int_{-2}^0f(x)\mathrm{\,d}x+\int_0^1f(x)\mathrm{\,d}x
			=\int_{-2}^0(x^2-2x+2)\mathrm{\,d}x+\int_0^1(\mathrm{e}^x+1)\mathrm{\,d}x  \\
			&=&\left(\dfrac{1}{3}x^3-x^2+2x\right)\bigg|_{-2}^0+ (\mathrm{e}^x+x)\bigg|_0^1=\dfrac{32}{3}+\mathrm{e}. 
		\end{eqnarray*}
		Vậy $\heva{&a=32 \\&b=3  \\&c=1}\Rightarrow a+b+c=36$.
	 }
\end{bt}

% nguồn: Phan2_Lop_12_Dai_so_Chuong_4
\begin{bt}%[2D4H2-5]%[ma:1ad5caed]
	Cho  hàm số $y=f(x)$ liên tục trên $\mathbb{R}$ thỏa mãn $\displaystyle\int_{-\dfrac{\pi}{4}}^{\frac{\pi}{4}}\left[\cos 2x+2f(x)\right]\mathrm{\,d}x=5$. Tính $\displaystyle\int_{-\frac{\pi}{4}}^{\frac{\pi}{4}}f(x)\mathrm{\,d}x$.
	 \shortans{2}
	\loigiai{ 
		Ta có \begin{align*}
		 5&=\displaystyle\int_{-\frac{\pi}{4}}^{\frac{\pi}{4}}[\cos 2x+2f(x)] \mathrm{\,d}x\\
		&=\int_{-\frac{\pi}{4}}^{\frac{\pi}{4}}\cos 2x \mathrm{\,d}x+\displaystyle\int_{-\frac{\pi}{4}}^{\frac{\pi}{4}}2f(x)\mathrm{\,d}x\\ 
		&=\frac{1}{2}\sin 2x\Big|_{\frac{-\pi}{4}}^{\frac{\pi}{4}}+2\displaystyle\int_{\frac{-\pi}{4}}^{\frac{\pi}{4}}f(x)\mathrm{\,d}x\\ 
		&=\frac{1}{2}\Big[\sin\dfrac{\pi}{2}-\sin\Big(-\dfrac{\pi}{2}\Big)\Big]+2\displaystyle\int_{\frac{\pi}{4}}^{\frac{\pi}{4}}f(x)\mathrm{\,d}x\\ 
		&=1+2\int_{\frac{\pi}{4}}^{\frac{\pi}{4}}f(x)\mathrm{\,d}x 
	\end{align*}
	Suy ra $\displaystyle\int_{-\frac{\pi}{4}}^{\frac{\pi}{4}}f(x)\mathrm{\,d}x=\dfrac{5-1}{2}=2$. 
} 
\end{bt}
